\subsection{The exact finite-level elliptic-unit comparison}
\label{subsec:cm-unit-index}

\begin{lemma}\label{lem:cm-unit-radical-comparison}
For every \(n\geq1\), the exact group \(C(L_n)\) in
\cite[Definition 2.2, p.~8]{IshiiSoule} has finite index in
\(O_{L_n}^{\times}\).
Consequently the two radical fields coincide:
\[
                         L_{\mathrm{ell}}=L_{\mathrm u}.
\]
\end{lemma}

\begin{proof}
We prove full rank with the defining theta values.
The formula below identifies their generators with the
Robert--Ramachandra invariants used in the logarithmic calculation.
Let \(F=L_n\), \(D=\operatorname{Gal}(F/K)\), and let
\(\chi\ne1\) be a complex character of \(D\).
Its conductor \(\mathfrak m\mid(p^n)\) is nontrivial, because
\(K\) has class number one.
Put \(F_{\mathfrak m}=K(\mathfrak m)\), and write
\(\chi_{\mathfrak m}\) for the induced primitive character.

Here is the required exact normalization.
For a lattice \(L\), the function
\[
 \varphi(z,L)=
       \bigl(e^{-\eta(z,L)z/2}\sigma(z,L)\bigr)^{12}\Delta(L)
\]
in \cite[(2), p.~255]{Oukhaba2009} is the fundamental theta
function in \cite[Definition 3.1(i), p.~6]{GotoIshii2026}.
The quotient formula in
\cite[Definition 3.1(ii), equation (1), p.~6]{GotoIshii2026}
gives, with precisely the rational function used by Ishii,
\[
 \theta_{\mathfrak a}(z;L)=
              \frac{\varphi(z,L)^{N\mathfrak a}}
                   {\varphi(z,\mathfrak a^{-1}L)}.
\]
Write \(L=\Omega_E O_K\), \(\mathfrak m=(\beta)\), and
\(\tau=\Omega_E/\beta\bmod L\).
Scaling both \(z\) and \(L\) leaves \(\varphi\) unchanged:
the factor \(\sigma^{12}\) scales by the twelfth power,
\(\Delta\) by its inverse, and \(\eta(z,L)z\) is invariant.
Thus \((\tau,L)\) may be replaced by \((1,\mathfrak m)\).

Let \(e_{\mathfrak m}>0\) generate
\(\mathfrak m\cap\mathbb Z\).
Oukhaba's invariants and reciprocity, equations (3)--(4) on
p.~255, state
\[
 \varphi_{\mathfrak m}(\sigma_{\mathfrak a})
      =\varphi(1,\mathfrak a^{-1}\mathfrak m)^{e_{\mathfrak m}},
 \qquad
 \varphi_{\mathfrak m}(1)^{\sigma_{\mathfrak a}}
      =\varphi_{\mathfrak m}(\sigma_{\mathfrak a}).
\]
Combining the displayed formulas gives the exact algebraic identity
\begin{equation}\label{eq:cm-exact-unit-convention}
 \theta_{\mathfrak a}(\tau)^{e_{\mathfrak m}}
       =\varphi_{\mathfrak m}(1)^{
                                  N\mathfrak a-\sigma_{\mathfrak a}}.
\end{equation}
The integer \(e_{\mathfrak m}\), which may be divisible by \(p\),
has been retained.

For \(u\in F^\times\), define
\[
 \ell_\chi^F(u)=
       \sum_{\sigma\in D}\chi(\sigma)^{-1}
                                    \log|\sigma(u)|^2.
\]
Then
\(\ell_\chi^F(u^\delta)=\chi(\delta)\ell_\chi^F(u)\).
For \(u\in F_{\mathfrak m}^\times\), one has
\[
 \ell_\chi^F(u)=[F:F_{\mathfrak m}]
                       \ell_{\chi_{\mathfrak m}}^{F_{\mathfrak m}}(u).
\]
The Kronecker limit formula
\cite[(6), p.~256]{Oukhaba2009}, applied at the primitive
conductor \(\mathfrak m\), gives
\[
 \ell_{\chi_{\mathfrak m}}^{F_{\mathfrak m}}
                         (\varphi_{\mathfrak m}(1))
   =-12r_{\mathfrak m}e_{\mathfrak m}
                          L'(0,\chi_{\mathfrak m}^{-1})\ne0 .
\]
Here \(r_{\mathfrak m}\) counts roots of unity of \(K\)
congruent to \(1\bmod\mathfrak m\).
The simple zero and nonzero derivative for a nontrivial primitive
character are also recorded explicitly in
\cite[\S2.3, p.~6, before (2.3)--(2.4)]{Oukhaba2003}.

Take \(c_0=1+6p\) and \(\mathfrak a=(c_0)\).
This auxiliary ideal is prime to \(6\mathfrak m\), and
\(N\mathfrak a-\chi_{\mathfrak m}(\sigma_{\mathfrak a})\ne0\):
the first term is the real integer \(c_0^2>1\) and the second
is a root of unity.
Equation \eqref{eq:cm-exact-unit-convention} therefore gives
\[
 \ell_\chi^F(\theta_{\mathfrak a}(\tau))
  =-12r_{\mathfrak m}[F:F_{\mathfrak m}]
    \bigl(N\mathfrak a-\chi_{\mathfrak m}(\sigma_{\mathfrak a})\bigr)
                         L'(0,\chi_{\mathfrak m}^{-1})\ne0.
\]
Choose \(\delta\in D\) with \(\chi(\delta)\ne1\).
The element
\[
                       v_\chi=\theta_{\mathfrak a}(\tau)^{\delta-1}
\]
is a generator allowed in Ishii's exact definition of \(C(F)\).
Indeed, \(F_{\mathfrak m}\subseteq F\), so the norm in that
definition is the identity.
The unit-injectivity requirement for \(\mathfrak m\) holds:
the norm of a difference of distinct roots of unity in an
imaginary quadratic field has prime divisors only \(2\) and \(3\),
while \(\mathfrak m\) is supported above \(p\geq5\).
For one-prime conductor the augmentation is still a global unit,
as in that definition; this also follows from the equal valuation
at all primes over that conductor in
\cite[(5), p.~256]{Oukhaba2009}.
Its \(\chi\)-logarithm is
\[
                \ell_\chi^F(v_\chi)
                   =(\chi(\delta)-1)
                      \ell_\chi^F(\theta_{\mathfrak a}(\tau))\ne0 .
\]

Dirichlet's logarithmic embedding of
\(O_F^\times/\mu(F)\) spans the real augmentation space of
\(\mathbb R[D]\).
This follows by choosing the embeddings extending the fixed
complex embedding of \(K\), one from each complex place.
After complexification, each nontrivial character of \(D\)
occurs exactly once.
The image of \(C(F)\) is \(D\)-stable, and the preceding
construction gives its nonzero projection to every such line.
It therefore has the full unit rank.
Since it is a subgroup of the finitely generated unit group
and already contains its torsion, its index is finite.

Finally, for any \(u\in O_{L_n}^\times\) there is an integer
\(d>0\) with \(u^d\in C(L_n)\).
If \(g\in H\) fixes \(L_{\mathrm{ell}}\), its exact Kummer
character satisfies \(d\,k_u(g)=0\) in \(\mathbb Z_p\).
Torsionfreeness gives \(k_u(g)=0\).
Thus \(g\) fixes all \(p\)-power roots of every global unit,
even when \(p\mid d\).
The reverse inclusion of radical fields follows from
\(C(L_n)\subseteq O_{L_n}^\times\), proving
\(L_{\mathrm{ell}}=L_{\mathrm u}\).
\end{proof}

\begin{theorem}\label{thm:cm-finite-plus-abelian}
Let \(\mathcal O\) be a \(c\)-stable set of characters of
\(\Delta=\operatorname{Gal}(K(p)/K)\), and suppose that
\((e_{\mathcal O}\mathcal X_2)^+\) is finite.
With
\(\mathcal O^\dagger=\{\omega_{\mathrm{cyc}}\chi^{-1}:
                                         \chi\in\mathcal O\}\),
the canonical field
\[
 L_{\mathcal O}=
 \Omega^{\,\ker(H\longrightarrow A
                \xrightarrow{\,e_{\mathcal O^\dagger}\,}
                              e_{\mathcal O^\dagger}A)}
\]
is contained in \(\Omega^*\).
Equivalently the image in \(A\) of the outer kernel has
zero projection to \(e_{\mathcal O^\dagger}A\).

For the set of all characters, finiteness of
\(\mathcal X_2^+\) gives
\[
 \Omega^{[H,H]}=L_{\mathrm{ell}}=L_{\mathrm u}
                   \subseteq\Omega^*,\qquad
 \bigcap_{\boldsymbol m\in I}\ker(\kappa_{\boldsymbol m}|_A)=0.
\]
The source \(H\) is then free pro-\(p\), by
Proposition~\ref{prop:cm-finite-plus-free}.
\end{theorem}

\begin{proof}
Proposition~\ref{prop:cm-component-unit-detection} gives
\(L_{\mathcal O}\subseteq L_{\mathrm u}\).
The preceding lemma identifies
\(L_{\mathrm u}=L_{\mathrm{ell}}\), and
Proposition~\ref{prop:cm-full-radical} places this field in
\(\Omega^*\).
These are actual inclusions inside \(\Omega\), so their
composition proves the component assertion.
For all characters, Proposition~\ref{prop:cm-finite-plus-free}
identifies \(L_{\mathrm u}\) with the maximal abelian
subextension \(\Omega^{[H,H]}\).
Corollary~\ref{cor:cm-radical-joint-kernel} then gives the
stated joint-character injectivity.
\end{proof}
